\section*{第 \ref{ch:b3:lp} 章 --- $L^p$ 空间}

\begin{solution}{exo:b3:lp:1}
(a) 在\cref{thm:b3:lp:holder}中取 $p = q = 2$：$\abs{\int f\bar
g\,\dd\mu} \leq \int\abs{fg} \leq \norm f_2\norm g_2$ ---
柯西--施瓦茨不等式，等号成立当且仅当 $\abs f, \abs g$
成比例，且相位一致。

(b) 在概率空间上，对 $p \leq q$：应用琴生不等式
（\cref{exo:b3:lp:10}），取凸函数 $\Phi(t) =
\abs t^{q/p}$，作用于函数 $\abs f^p$：
$\bigl(\int\abs f^p\bigr)^{q/p} \leq \int\abs f^q$，即
$\norm f_p \leq \norm f_q$。

(c) $\int\abs{fg} = \norm f_1\norm g_\infty$ 当且仅当 $\abs{g} =
\norm g_\infty$ 在 $\{f \neq 0\}$ 上几乎处处成立（不等式
$\abs{fg} \leq \abs f\norm g_\infty$ 须几乎处处
取等）。
\end{solution}

\begin{solution}{exo:b3:lp:2}
$\int_0^1 x^{-p/2}\dd x < \infty$ 当且仅当 $p < 2$：第一项
属于 $L^p$，当 $p \in \intco12$ 时。$\int_1^\infty x^{-p/2}\dd x <
\infty$ 当且仅当 $p > 2$：第二项当 $p \in \intoo2\infty$ 时（并
$p = \infty$ 时它有界 --- 应一并计入）。第三项当 $p <
2$ 时，由 $x^{-p/2}$ 控制，可积；对 $p = 2$，
作代换 $u = -\ln x$：$\int_0^1\frac{\dd x}{x(1 + \abs{\ln
x})^2} = \int_0^\infty\frac{\dd u}{(1 + u)^2} < \infty$；当
$p > 2$ 时，幂次占优，发散。故 $p \in \intcc12$。
第四项：$\int_\R\frac{\dd x}{(1 + \abs x)^p} < \infty$ 当且仅当 $p >
1$；有界，因而也属于 $L^\infty$：$p \in
\intoc1\infty$。由此可见：在 $0$ 附近，可积性偏好较小的 $p$；在
$\infty$ 附近偏好较大的 $p$。综合这两点，
$L^p(\R)$ 空间之间不存在包含关系。
\end{solution}

\begin{solution}{exo:b3:lp:3}
(a) $\norm{f_n}_p^p = \lambda(I_n) \to 0$（第 $k$ 层上的
区间长度为 $2^{-k}$）。但在每一层中，每个 $x$ 恰落在一个区间内：
$f_n(x) = 1$ 无穷多次，且 $= 0$
无穷多次（同一层中不含
$x$ 的区间）。故在任何点处都不收敛。

(b) $f_{n_k} = \mathbf 1_{\intcc0{2^{-k}}}$（每层的第一个
区间）收敛于 $0$，对一切 $x > 0$ 成立，故几乎处处收敛。

(c) 几乎处处收敛，但不属于 $L^1$：$n\mathbf 1_{\intoo0{1/n}} \to 0$
几乎处处，积分为 $1$。属于 $L^1$ 但不属于任何 $L^p$（$p > 1$）：$g_n =
\eu^n\,\mathbf 1_{(0,\ \eu^{-n}/n)}$：$\norm{g_n}_1 = \frac1n
\to 0$，而 $\norm{g_n}_p^p = \eu^{(p-1)n}/n \to \infty$ 对
每个 $p > 1$ 成立。
\end{solution}

\begin{solution}{exo:b3:lp:4}
上界：$\norm f_p \leq \mu(X)^{1/p}\norm f_\infty$
（由 \cref{prop:b3:lp:inclusions}(a) 取 $q = \infty$），且
$\mu(X)^{1/p} \to 1$。下界：对 $\varepsilon > 0$，$A =
\{\abs f > \norm f_\infty - \varepsilon\}$ 满足 $\mu(A) > 0$
（由本质上确界的定义），且
\[
\norm f_p \geq \Bigl(\int_A \abs f^p\Bigr)^{1/p}
\geq (\norm f_\infty - \varepsilon)\,\mu(A)^{1/p}
\xrightarrow[p\to\infty]{} \norm f_\infty - \varepsilon .
\]
\end{solution}

\begin{solution}{exo:b3:lp:5}
(a) 有两点：转换 $\norm{\tau_hg - g}_p \leq
C^{1/p}\norm{\tau_hg - g}_\infty$（支集测度有限）
在 $p = \infty$ 时失效，仅在于 \emph{
$\mathcal C_c$ 的稠密性}在该处不成立 --- 这才是真正的缺口：
\cref{thm:b3:lp:density} 的第 (2) 步没有 $L^\infty$ 上的对应结论。

(b) 若 $f$ 有一致连续的代表元 $g$：
$\norm{\tau_hf - f}_\infty = \sup_x\abs{g(x - h) - g(x)} \to
0$。反之，设 $\norm{\tau_hf - f}_\infty \to 0$。
光滑化 $f_\varepsilon = f * \rho_\varepsilon$ 是
连续的，且
\[
\norm{f_\varepsilon - f}_\infty
\leq \sup_{\norm y \leq \varepsilon}\norm{\tau_yf -
f}_\infty \longrightarrow 0
\]
（卷积是若干平移的平均）。每个
$f_\varepsilon$ 一致连续
（$\norm{\tau_hf_\varepsilon - f_\varepsilon}_\infty \leq
\norm{\tau_hf - f}_\infty$，由平均得），且一致连续函数的
一致极限仍一致连续：$f$ 与
某个一致连续函数几乎处处相等。
\end{solution}

\begin{solution}{exo:b3:lp:6}
由赫尔德不等式，对每个 $x$，被积函数 $y \mapsto f(x -
y)g(y)$ 属于 $L^1$，且 $\abs{(f*g)(x)} \leq \norm f_p\norm
g_q$：处处有定义且有界。一致连续性（$p <
\infty$）：
\[
\abs{(f*g)(x + h) - (f*g)(x)}
= \Bigl|\int\bigl(\tau_{-h}f - f\bigr)(x - y)\,g(y)\dd y\Bigr|
\leq \norm{\tau_{-h}f - f}_p\,\norm g_q
\xrightarrow[h\to0]{} 0,
\]
关于 $x$ 一致（\cref{thm:b3:lp:density}(3)）。若 $p =
\infty$，则 $q = 1$：写 $f * g = g * f$，并用同样的
估计，此时平移作用在 $g \in L^1$ 上。
\end{solution}

\begin{solution}{exo:b3:lp:7}
取定 $\chi \in \mathcal C_c^\infty(\intoo ab)$，满足 $\int\chi =
1$，并令 $c = \int f\chi$。任取 $\varphi \in \mathcal
C_c^\infty$，并令 $\psi = \varphi -
\bigl(\int\varphi\bigr)\chi$：则 $\int\psi = 0$，故 $\Phi(x)
= \int_a^x\psi$ 定义了 $\Phi \in \mathcal C_c^\infty(\intoo
ab)$（它在两端附近为零：在 $a$ 附近显然，在 $b$
附近因总积分为 $0$），且 $\Phi' = \psi$。
假设给出 $\int f\psi = \int f\Phi' = 0$，从而
\[
\int f\varphi = \Bigl(\int\varphi\Bigr)\int f\chi = \int
c\,\varphi
\quad\text{对每个 } \varphi:
\qquad \int(f - c)\varphi = 0 .
\]
由 \cref{cor:b3:lp:fundlemma}（局部化到 $\intoo ab$ 上），$f =
c$ 几乎处处成立。
\end{solution}

\begin{solution}{exo:b3:lp:8}
设 $3\delta < d(K, \R^d\setminus U)$（其为正，
见\cref{exo:b3:topology:6}(b)；若 $U = \R^d$，则任意 $\delta$
均可），记 $K_\delta = \{x : d(x, K) \leq \delta\}$，且
\[
\varphi = \mathbf 1_{K_\delta} * \rho_{\delta/2} .
\]
于是 $\varphi \in \mathcal C^\infty$
（见\cref{thm:b3:lp:regularization}(1)；示性函数属于 $L^1$），
$0 \leq \varphi \leq 1$（$\int\rho = 1$），$\varphi = 1$ 在 $K$ 上，
（对 $x \in K$，$\bar B(x, \delta/2) \subseteq K_\delta$，故
卷积完整地积分了 $\rho$），且
$\operatorname{supp}\varphi \subseteq K_{3\delta/2} \subseteq
U$：紧支集（$K_\delta$ 有界）。关于
单位分解：给定 $K \subseteq U_1\cup\dots\cup U_m$，由
紧性取紧集 $K_i \subseteq U_i$，使 $K \subseteq
\bigcup \mathring K_i$，取如上的 $\varphi_i$，对 $(K_i,
U_i)$，并令 $\psi_i = \varphi_i\prod_{j <i}(1 -
\varphi_j)$：每个 $\psi_i \in \mathcal C_c^\infty(U_i)$，且
$\sum_i\psi_i = 1 - \prod_i(1 - \varphi_i) = 1$ 在 $K$ 上。
\end{solution}

\begin{solution}{exo:b3:lp:9}
(a) 对 $(p_0, q_0) = (1, \infty)$，
在 $r = p$ 处的插值指数为 $\alpha = \frac1p$：
由\cref{prop:b3:lp:inclusions}(c)得 $\norm f_p \leq \norm
f_1^{1/p}\norm f_\infty^{1 - 1/p}$。

(b) 在
\cref{prop:b3:lp:inclusions}(c)中取对数，得 $\ln\norm f_r \leq
\alpha\ln\norm f_p + (1-\alpha)\ln\norm f_q$，其中 $\frac1r$
是 $\frac1p, \frac1q$ 的同一凸组合：$\frac1p
\mapsto \ln\norm f_p$ 是凸的。取一个例子，其 $L^p$ 成员资格
恰好在 $\intoo{p_0}{p_1}$ 上成立：
\[
f(x) = x^{-1/p_1}\,\mathbf 1_{\intoo01}(x) +
x^{-1/p_0}\,\mathbf 1_{\intco1\infty}(x):
\]
第一项属于 $L^p$ 当且仅当 $p < p_1$，第二项当且仅当 $p >
p_0$。
\end{solution}

\begin{solution}{exo:b3:lp:10}
设 $m = \int f\,\dd\mu \in \R$。凸性给出一条
在 $m$ 处的支撑直线：存在 $s$，使得 $\Phi(t) \geq \Phi(m) + s(t -
m)$ 对所有 $t$ 成立（取 $s$ 在单侧导数之间，
凸函数的单侧导数总是存在的）。取 $t = f(x)$ 代入，
并关于概率测度 $\mu$ 积分：
\[
\int\Phi\circ f\,\dd\mu \geq \Phi(m) + s\Bigl(\int f - m\Bigr)
= \Phi\Bigl(\int f\,\dd\mu\Bigr)
\]
（可测性：$\Phi$ 连续；
$\Phi\circ f$ 负部的可积性由支撑直线保证）。
算术–几何平均不等式：在有限集上权为 $w_i$，
取 $\Phi = \exp$，$f = \sum(\ln a_i)\mathbf 1_i$：
$\exp\bigl(\sum w_i\ln a_i\bigr) \leq \sum w_ia_i$，即
$\prod a_i^{w_i} \leq \sum w_ia_i$。范数的单调性见
\cref{exo:b3:lp:1}(b)。
\end{solution}

\begin{solution}{exo:b3:lp:11}
(a) 先假设 $p, q, r < \infty$ 且 $f, g \geq 0$
（可换为绝对值）。三个指数 $r$，
$\alpha = \frac{pr}{r-p}$, $\beta = \frac{qr}{r-q}$
满足 $\frac1r + \frac1\alpha + \frac1\beta = \frac1r +
\frac1p - \frac1r + \frac1q - \frac1r = 1$（标度
关系）。对固定的 $x$，作分解
\[
f(y)g(x-y) = \bigl[f(y)^pg(x-y)^q\bigr]^{1/r}\cdot
f(y)^{1 - p/r}\cdot g(x-y)^{1 - q/r},
\]
由三个指数的赫尔德不等式得
\[
(f*g)(x) \leq \Bigl(\int f^pg(x-\cdot)^q\Bigr)^{1/r}
\norm f_p^{\,p/\alpha\cdot\alpha/p}\cdots
\]
更准确地说，第二个因子为
$\bigl(\int f^{(1-p/r)\alpha}\bigr)^{1/\alpha} =
\norm f_p^{p(1/p - 1/r)}$，因为 $(1 - \frac pr)\alpha = p$，
同理第三个因子为 $\norm g_q^{q(1/q - 1/r)}$。两边
取 $r$ 次幂，并对 $x$ 积分（对第一个因子用
托内利定理）：
\[
\norm{f*g}_r^r \leq \norm f_p^p\,\norm g_q^q\cdot
\norm f_p^{\,rp(1/p - 1/r)}\,\norm g_q^{\,rq(1/q - 1/r)}
= \norm f_p^r\,\norm g_q^r .
\]
端点情形（$r = \infty$，或某指数等于其
界）则为普通的赫尔德不等式或直接估计。

(b) $r = \infty$ 迫使 $q = p'$：$\abs{f*g(x)} \leq \norm
f_p\norm g_{p'}$ --- 经平移与反射后用赫尔德不等式。
$q = 1$ 时 $r = p$：$\norm{f*g}_p \leq \norm
g_1\norm f_p$，这是磨光化的主力工具
（\cref{thm:b3:lp:regularization}的核心）。$p = q = 1$
时 $r = 1$：卷积代数
（见\cref{thm:b3:product:convolution}）。

(c) 将 $f, g$ 换为 $f_\lambda = f(\lambda\cdot)$，
$g_\lambda = g(\lambda\cdot)$，则 $f_\lambda*g_\lambda =
\lambda^{-d}(f*g)(\lambda\cdot)$，比较范数得
\[
\text{左边} \sim \lambda^{-d - d/r}, \qquad
\text{右边} \sim \lambda^{-d/p - d/q} :
\]
对所有 $f, g$ 成立的不等式要求两个标度指数
相等，即 $1 + \frac1r = \frac1p + \frac1q$
恰好成立。其他任何组合都会在 $\lambda \to 0$ 或
$\infty$ 时失效：杨氏关系式并非权宜之计，而是
标度律。
\end{solution}

\begin{solution}{exo:b3:lp:12}
(a) 不妨设 $\norm f_p = \norm g_q = 1$。赫尔德不等式的证明
是对杨氏不等式 $\abs{fg} \leq
\frac{\abs f^p}p + \frac{\abs g^q}q$ 作积分，积分取等
迫使杨氏不等式几乎处处取等，而由
$\exp$ 的严格凸性（取等当且仅当 $a^p = b^q$）得 $\abs
f^p = \abs g^q$ 几乎处处成立。撤去归一化：
$\abs f^p\norm g_q^q = \abs g^q\norm f_p^p$ 几乎处处成立 ---
即成比例。

(b) 闵可夫斯基不等式即对 $\abs{f + g}^{p-1}
\abs f$ 与 $\abs{f+g}^{p-1}\abs g$ 分别应用两次赫尔德不等式；取等时，(a) 中的
成比例关系在两处均成立：$\abs f^p$ 与 $\abs g^p$ 各自
与 $\abs{f+g}^{(p-1)q} = \abs{f+g}^p$ 成比例，因此
$\abs g = t\abs f$ 几乎处处成立，其中常数 $t \geq 0$；而
逐点三角不等式 $\abs{f + g} \leq
\abs f + \abs g$ 也必须几乎处处取等，对复值而言，这意味着
$f$ 与 $g$ 在二者均非零处几乎处处具有相同的
幅角。综合得 $g = tf$ 几乎处处成立，
$t > 0$（二者均非零）。

(c) 对 $p = 1$，等式 $\int\abs{f + g} = \int\abs f +
\int\abs g$ 在 $f, g$ 有相同符号模式时总成立
（几乎处处幅角相同） --- 无需成比例：
$f = \mathbf 1_{\intcc01}$ 与 $g = \mathbf 1_{\intcc02}$
均满足要求。对 $p = \infty$，$\norm{f+g}_\infty = \norm f_\infty +
\norm g_\infty$ 成立，只要两函数
相容地在某一公共点（或沿某公共点列）上取到峰值：
$f = g\,$ 在某一点附近成立即可，无论其他处
的行为如何。$L^p$ 球的严格凸性在 $1 <
p < \infty$ 时成立 --- 以及其在端点处的失效 ---
正是这些取等情形所体现的。
\end{solution}

\begin{solution}{pb:b3:lp:1}
\textbf{1.} $f$ 的支集位于某个区间 $[\alpha, \beta]
\subseteq \intoo0\infty$ 内，于是 $F = 0$ 在 $[0, \alpha]$ 上，且 $F
\equiv F(\beta)$ 在 $[\beta, \infty)$ 上，$Hf$ 在 $0$ 附近为零，
且在无穷远处为 $O(1/x)$；$\int_1^\infty x^{-p}\dd x <
\infty$（$p > 1$），且 $Hf$ 连续，故 $Hf \in L^p$。

\textbf{2.} $\bigl(x^{1-p}F(x)^p\bigr)' = (1-p)x^{-p}F^p +
p\,x^{1-p}F^{p-1}f$。两端的边界值均为零：在 $0$ 处
因 $F = 0$ 在 $0$ 附近成立；在 $\infty$ 处因 $x^{1-p}F^p \leq
F(\beta)^p x^{1-p} \to 0$（$p > 1$）。对该恒等式积分
在 $\intoo0\infty$ 上：
\[
0 = (1 - p)\int_0^\infty\Bigl(\frac Fx\Bigr)^p\dd x
+ p\int_0^\infty\Bigl(\frac Fx\Bigr)^{p-1}f(x)\,\dd x,
\]
即得所示关系式。

\textbf{3.} 取指数为 $q = \frac p{p-1}$ 与 $p$ 的赫尔德不等式：
\[
\int\Bigl(\frac Fx\Bigr)^{p-1}f
\leq \Bigl(\int\Bigl(\frac
Fx\Bigr)^{p}\Bigr)^{1 - 1/p}\,\norm f_p ,
\]
因此 $\norm{Hf}_p^p \leq \frac p{p-1}\norm{Hf}_p^{p-1}\norm
f_p$，两边作除法（由第 1 问知其有限；若为 $0$，则
无需证明）。

\textbf{4.} 设 $f \in L^p$，$f \geq 0$。取 $\varphi_k \in
\mathcal C_c(\intoo0\infty)$ 使 $\varphi_k \to f$ 于 $L^p$ 中
（由\cref{thm:b3:lp:density}(2)，与开
半直线相交 --- 逼近 $f\mathbf 1_{[1/k, k]}$，并
用对角线法）；再将 $\varphi_k$ 换为 $\abs{\varphi_k}$
（仍连续，且更接近 $f \geq 0$）。对每个固定的
$x > 0$，在 $\intoo0x$ 上由赫尔德不等式得
\[
\abs{H\varphi_k(x) - Hf(x)} \leq
\frac1x\,x^{1 - 1/p}\,\norm{\varphi_k - f}_p \to 0 :
\]
$H\varphi_k \to Hf$ 逐点成立。由法图引理及第 3 问：
\[
\int (Hf)^p \leq \liminf_k\int(H\varphi_k)^p
\leq \Bigl(\frac p{p-1}\Bigr)^p\liminf_k\norm{\varphi_k}_p^p
= \Bigl(\frac p{p-1}\Bigr)^p\norm f_p^p .
\]
对变号或复值的 $f$：逐点有 $\abs{Hf} \leq H\abs f$，
故非负情形可用于 $\abs f$。

\textbf{5.} $\norm{f_A}_p^p = \int_1^A t^{-1}\dd t = \ln A$.
对 $1 \leq x \leq A$：
\[
(Hf_A)(x) = \frac1x\int_1^x t^{-1/p}\dd t
= \frac{x^{1 - 1/p} - 1}{(1 - \tfrac1p)\,x}
= \frac p{p-1}\,x^{-1/p}\bigl(1 - x^{-(1 - 1/p)}\bigr).
\]

\textbf{6.} 固定 $\varepsilon > 0$，取 $X_0$ 使得 $(1 -
x^{-(1-1/p)})^p \geq 1 - \varepsilon$，对 $x \geq X_0$ 成立。于是
\[
\norm{Hf_A}_p^p \geq \int_{X_0}^A\Bigl(\frac
p{p-1}\Bigr)^p\frac{1 - \varepsilon}{x}\,\dd x
= \Bigl(\frac p{p-1}\Bigr)^p(1 - \varepsilon)\,(\ln A - \ln
X_0),
\]
因此 $\dfrac{\norm{Hf_A}_p^p}{\norm{f_A}_p^p} \geq \bigl(\frac
p{p-1}\bigr)^p(1 - \varepsilon)\bigl(1 - \frac{\ln X_0}{\ln
A}\bigr) \to \bigl(\frac p{p-1}\bigr)^p(1 - \varepsilon)$ 当
$A \to \infty$ 时，常数不可改进。

\textbf{7.} 第 3 问中的取等迫使
赫尔德不等式取等：$f^p$ 与 $\bigl(\frac Fx\bigr)^{(p-1)q}
= \bigl(\frac Fx\bigr)^p$ 几乎处处成立，即 $f = \gamma\,\frac Fx$
几乎处处成立，其中 $\gamma \geq 0$ 为某常数。由于 $F(x) = \int_0^xf$
绝对连续，且 $F' = f$ 几乎处处成立，$F$ 满足 $F' =
\gamma F/x$：在任一使 $F > 0$ 的区间上，$(\ln F)' =
\gamma/x$，因此 $F = c\,x^{\gamma}$，且 $f = c\gamma
x^{\gamma - 1}$ 在该区间上。但没有非零幂 $x^{\gamma-1}$
属于 $L^p(\intoo0\infty)$（需 $p(\gamma - 1) < -1$
在 $\infty$ 处成立，而 $> -1$ 为 $0$ 处所需：二者不相容），又 $F$
不能恒为零，除非 $f = 0$。因此取等要求 $f =
0$。

\textbf{8.} 设 $(a_n)$ 非增且 $\geq 0$，定义
阶梯函数 $g(t) = a_{\lceil t\rceil}$，定义在
$\intoo0{+\infty}$ 上非增，且 $\int_0^\infty g^p =
\sum_na_n^p$，且 $\int_0^ng = a_1 + \dots + a_n$，因此
$(Hg)(n) = \frac{a_1 + \dots + a_n}n$。非增函数的平均
仍非增，因此 $Hg$ 亦然，于是
\[
\sum_{n\geq1}\Bigl(\frac{a_1{+}\dots{+}a_n}n\Bigr)^p
= \sum_{n\geq1}(Hg)(n)^p
\leq \sum_{n\geq1}\int_{n-1}^n (Hg)(t)^p\,\dd t
= \norm{Hg}_p^p
\leq \Bigl(\frac p{p-1}\Bigr)^p\sum_n a_n^p
\]
由第一部分可得。对一般的非负数列，设 $(a_n^*)$
为其非增重排（这在 $a_n \to 0$ 时可以做到，
这一条件可以假设 --- 否则两边都是无穷大）：
右边保持不变，而每个部分和 $a_1 + \dots
+ a_n$ 不超过 $a_1^* + \dots + a_n^*$（即前 $n$ 个最大项之和），
左边只增不减。于是对一切
$(a_n)$ 不等式成立。

\textbf{9.} 若 $(a_n) \in \ell^p$，则切萨罗
平均数列属于 $\ell^p$，且范数 $\leq \frac p{p-1}\norm
a_p$。例（$p = 2$）：$a_n = \frac1{\sqrt n\,\ln n}$（$n
\geq 2$）：$\sum a_n^2 = \sum\frac1{n\ln^2n} < \infty$，然而
$\sum a_n = \infty$（积分判别法）；哈代不等式仍保证
$\sum_n\bigl(\frac{a_1 + \dots + a_n}n\bigr)^2 < \infty$。

\textbf{10.} 对 $f = \mathbf 1_{\intcc01}$：$Hf(x) = 1$ 在
$\intoc01$ 上，且 $= \frac1x$（$x \geq 1$）：$\int_0^\infty Hf =
1 + \int_1^\infty\frac{\dd x}x = \infty$，而 $\norm f_1 =
1$。证明在两点处失效：常数
$\frac p{p-1}$ 在 $p \to 1$ 时趋于无穷，而边界项
$x^{1-p}F^p$ 在 $p = 1$ 时于无穷远处不再趋于零。
哈代不等式是真正的 $p > 1$ 现象。

\textbf{11.} 取最长的区间 $B_{i_1}$；丢弃
所有与之相交的区间；在剩下的区间中取最长者
$B_{i_2}$；如此迭代（区间个数有限）。所选区间
由构造两两不交；每个被丢弃的 $B$ 都与某个
长度不小于它的所选区间相交。与较长或等长区间相交的
区间包含于其三倍区间中，即 $B \subseteq
3B_{i_l}$。

\textbf{12.} $\{Mf > t\}$ 为开集：每个平均 $x \mapsto
\frac1{2r}\int_{x-r}^{x+r}\abs f$ 关于 $x$ 连续（由控制收敛定理），
而连续函数的上确界是下半连续的。
集合中每个 $x$ 都拥有 $B_x =
\intoo{x-r_x}{x+r_x}$，满足 $\int_{B_x}\abs f >
t\,\lambda(B_x)$。对紧集 $K \subseteq \{Mf > t\}$：
有限个 $B_x$ 覆盖 $K$，由第 11 题的维塔利引理取出
两两不交的 $B_1', \dots, B_k'$，使得 $K \subseteq
\bigcup_l3B_l'$，于是
\[
\lambda(K) \leq 3\sum_l\lambda(B_l') <
\frac3t\sum_l\int_{B_l'}\abs f \leq \frac3t\,\norm f_1
\]
由不交性可得。由内正则性
（\cref{thm:b3:measure:regularity}）即得结论。

\textbf{13.} 对 $x > 1$：当 $r \in \intcc{x-1}x$ 时
平均值为 $\frac{r - x + 1}{2r}$，关于 $r$ 递增；当 $r
\geq x$ 时为 $\frac1{2r}$，递减：上确界为
$\frac1{2x}$，在 $r = x$ 处取到。故 $M\mathbf 1_{\intcc01}
\notin L^1$。一般地，若 $\int_I\abs f = c > 0$ 在一个
有界区间 $I \subseteq \intcc{-C}C$ 上，则 $Mf(x) \geq
\frac{c}{2(\abs x + C)}$ 对一切 $x$ 成立：它永远不可积，
除非 $f = 0$ 几乎处处成立。

\textbf{14.} 记 $f_t = f\,\mathbf 1_{\abs f > t/2}$，则
$M(f - f_t) \leq \frac t2$，于是 $\{Mf > t\} \subseteq \{Mf_t
> \frac t2\}$，由第 12 题得 $\lambda(Mf > t) \leq
\frac6t\int_{\abs f > t/2}\abs f$。层饼表示
（\cref{prop:b3:product:layercake}）与托内利定理给出：
\[
\norm{Mf}_p^p = p\int_0^\infty t^{p-1}\lambda(Mf > t)\dd t
\leq 6p\int\abs f\int_0^{2\abs f}t^{p-2}\,\dd t\,\dd\lambda
= \frac{6p\,2^{p-1}}{p-1}\,\norm f_p^p .
\]

\textbf{15.} 记 $A_rf(x) = \frac1{2r}\int_{x-r}^{x+r}f$。
对连续的 $g$：$A_rg(x) \to g(x)$ 处处成立。给定
$\varepsilon > 0$，将 $f = g + h$ 分解，其中 $g$ 是
紧支集的连续函数，且 $\norm h_1 < \varepsilon$
（\cref{thm:b3:lp:density}）；于是
\[
\limsup_{r\to0}\,\abs{A_rf(x) - f(x)} \leq Mh(x) +
\abs{h(x)},
\]
故 $\Omega_\delta = \{\limsup_r\abs{A_rf - f} > \delta\}
\subseteq \{Mh > \tfrac\delta2\} \cup \{\abs h >
\tfrac\delta2\}$ 的测度 $\leq \frac{6\varepsilon}\delta
+ \frac{2\varepsilon}\delta$（由第 12 题与马尔可夫不等式）。
$\varepsilon$ 任意，故 $\lambda(\Omega_\delta) = 0$；
对 $\delta = \frac1k$ 取并：$A_rf \to f$ 几乎处处成立。

\textbf{16.} (a) 对每个 $q \in \Q$，将第 15 题用于
$\abs{f - q}$，得 $A_r\abs{f - q}(x) \to \abs{f(x) - q}$
几乎处处成立；在这些全测度集的交集上，选取
$q$，使 $\abs{f(x) - q} < \eta$：
$\limsup_rA_r\abs{f - f(x)}(x) \leq 2\eta$ 对每个 $\eta$ 成立：
几乎每个 $x$ 都是勒贝格点。(b) 在勒贝格
点处，
\[
\Bigl|\frac{F(x + h) - F(x)}h - f(x)\Bigr|
= \Bigl|\frac1h\int_x^{x+h}\bigl(f - f(x)\bigr)\Bigr|
\leq 2\,A_{\abs h}\abs{f - f(x)}(x) \to 0 :
\]
$F' = f$ 几乎处处成立\ --- $L^1$ 函数的原函数确实
可以求导回去；阶梯函数（\cref{pb:b3:measure:1}）
仅是\emph{逆方向}的反例。

\textbf{17.} 对 $\mathbf 1_{A\cap[-n,n]}$ 应用第 15 题，
并令 $n$ 增大：对几乎处处的 $x \in A$，密度
$\frac{\lambda(A\cap\intcc{x-r}{x+r})}{2r} \to 1$。施坦豪斯：
在密度点附近取 $r$，使密度 $> \frac34$；
对 $\abs t < \frac r2$，$A$ 与 $A + t$ 各占超过
$\frac32r$ 的某区间（长度 $\leq \frac52r$）中的部分，故二者
相交，从而 $\intoo{-r/2}{r/2} \subseteq A - A$。

\textbf{18.} 令 $G(x) = x^{\alpha+1-p}F(x)^p$，$G$ 在
$0$ 处为零（$F$ 在 $0$ 附近为零），在 $\infty$ 处也为零（$F$ 有界，
$\alpha + 1 - p < 0$），故 $\int_0^\infty G' = 0$，其中
\[
G' = (\alpha + 1 - p)\,x^{\alpha-p}F^p +
p\,x^{\alpha+1-p}F^{p-1}f
= (\alpha + 1 - p)\Bigl(\frac Fx\Bigr)^px^\alpha +
p\Bigl(\frac Fx\Bigr)^{p-1}f\,x^\alpha .
\]
于是 $\int(F/x)^px^\alpha = \frac{p}{p-1-\alpha}
\int(F/x)^{p-1}f\,x^\alpha$；对测度
$x^\alpha\dd x$ 用赫尔德不等式（指数 $\frac p{p-1}$ 与 $p$），
与第 3 题相同即可完成。临界情形 $\alpha = p - 1$，取 $f(t) =
\frac1t\mathbf 1_{\intcc1A}$，右端为 $\ln A$
而左端包含 $\int_1^A\frac{(\ln x)^p}x\dd x =
\frac{(\ln A)^{p+1}}{p+1}$，故没有常数能在 $A \to
\infty$ 时仍成立。

\textbf{19.} 在 $\{0 < t < x\}$ 上用托内利定理：
\[
\langle Hf, g\rangle =
\int_0^\infty\frac{g(x)}x\int_0^xf(t)\,\dd t\,\dd x
= \int_0^\infty f(t)\int_t^\infty\frac{g(x)}x\,\dd x\,\dd t
= \langle f, H^*g\rangle .
\]
由对偶性，$\norm{H^*f}_p = \sup\{\langle f, Hg\rangle : g
\geq 0, \norm g_q \leq 1\} \leq \norm f_p\cdot
\sup\norm{Hg}_q \leq \frac q{q-1}\norm f_p = p\,\norm f_p$，
在 $L^q$ 中用哈代不等式，其常数 $\frac q{q-1}$
等于 $p$。

\textbf{20.} 沿（零测）对角线分开。在 $\{y \leq
x\}$：
\[
\iint_{y\leq x}\frac{f(x)g(y)}{x}\,\dd y\,\dd x
= \int f(x)\,(Hg)(x)\,\dd x \leq \norm f_p\,\norm{Hg}_q
\leq p\,\norm f_p\norm g_q,
\]
由于 $L^q$ 中的哈代不等式常数为 $\frac q{q-1} =
p$。对称地，$\iint_{y>x} = \int g\,(Hf) \leq
\norm g_q\norm{Hf}_p \leq q\,\norm f_p\norm g_q$（$\frac
p{p-1} = q$）。总和为 $(p + q)\,\norm f_p\norm g_q$。

\textbf{21.} 取 $a_n = n^{-1/p}$，$n \leq N$，则右边
为 $\bigl(\frac p{p-1}\bigr)^p\sum_{n\leq N}\frac1n =
\bigl(\frac p{p-1}\bigr)^p\ln N + O(1)$。左端中，对 $n
\leq N$：$a_1 + \dots + a_n \geq \int_1^{n+1}t^{-1/p}\dd t =
\frac{p}{p-1}\bigl((n+1)^{1-1/p} - 1\bigr)$，于是第 $n$
个切萨罗平均 $\geq \frac p{p-1}n^{-1/p}(1 -
o(1))$ 对 $n$ 在任意范围 $n \geq n_0(\eta)$ 内一致成立；
取 $p$ 次幂并求和，左端为 $\geq
\bigl(\frac p{p-1}\bigr)^p(1 - \eta)\ln N + O_\eta(1)$。
两端相除并令 $N \to \infty$，再令 $\eta \to 0$：
没有小于 $\bigl(\frac p{p-1}\bigr)^p$ 的常数能成立。

\textbf{22.} $H$ 在 $L^p$ 上有界：累积平均
不会放大 $p$ 范数（常数 $\frac p{p-1}$）；切萨罗平均在
$\ell^p$ 上同理，只是离散化；$M$ 在 $L^p$ 上有界：
即使最优局部平均也保持可控（常数
$O(\frac1{p-1})$）。三者在 $p = 1$ 处均失效，原因只有一个：
把集中的单位质量取平均，会产生一个
$\frac1x$ 型尾部（第 10、13 题），而 $\frac1x$ 属于
每个 $L^p$（在无穷远处），唯 $L^1$ 除外。光滑化把
质量恰好扩散到可积性的调和边界。

\textbf{23.} 令 $b_n = a_n^{1/p}$，于是 $\sum b_n^p = \sum
a_n$。由第 8 题可得
\[
\sum_{n\geq1}\Bigl(\frac{b_1 + \dots + b_n}n\Bigr)^{p}
\leq \Bigl(\frac p{p-1}\Bigr)^{p}\sum_{n\geq1}a_n,
\]
并由算术–几何平均不等式，每一项有如下下界：
\[
\Bigl(\frac{b_1 + \dots + b_n}n\Bigr)^{p}
\geq \bigl(b_1\cdots b_n\bigr)^{p/n}
= \bigl(a_1\cdots a_n\bigr)^{1/n}.
\]
于是 $\sum(a_1\cdots a_n)^{1/n} \leq \bigl(\frac
p{p-1}\bigr)^p\sum a_n$ 对\emph{每个} $p > 1$ 成立。令 $m =
p - 1$，
\[
\Bigl(\frac p{p-1}\Bigr)^{p} =
\Bigl(1 + \frac1m\Bigr)^{m+1},
\]
它递减趋于 $\eu$，当 $m \to \infty$ 时（
即 $\eu$ 的经典单调上界数列）。取
$p$ 的下确界，得到常数为 $\eu$ 的卡莱曼不等式。

\textbf{24.} 取 $a_n = \frac1n$，则 $(a_1\cdots a_n)^{1/n} =
(n!)^{-1/n}$。由\cref{pb:b3:product:1}的夹逼可得
$\sqrt{2\pi n}\,(n/\eu)^n \leq n! \leq \sqrt{2\pi
n}\,(n/\eu)^n\eu^{1/(12n)}$，故
\[
(n!)^{1/n} = \frac n\eu\,(2\pi n)^{1/(2n)}
\eu^{O(1/n^2)}
= \frac n\eu\Bigl(1 +
O\Bigl(\frac{\ln n}{n}\Bigr)\Bigr),
\]
由于 $(2\pi n)^{1/(2n)} = \exp\bigl(\frac{\ln(2\pi
n)}{2n}\bigr)$。取倒数，$(n!)^{-1/n} = \frac\eu n(1 +
O(\frac{\ln n}n))$，并对 $n \leq N$ 求和：
\[
\sum_{n\leq N}(n!)^{-1/n}
= \eu\ln N + O(1),
\qquad
\eu\sum_{n\leq N}\frac1n = \eu\ln N + O(1)
\]
（误差级数 $\sum\frac{\ln n}{n^2}$ 收敛）。若
常数为 $c$ 的卡莱曼不等式成立，则必有 $\eu\ln N
+ O(1) \leq c\,(\ln N + O(1))$，于是 $c \geq \eu$
两端除以 $\ln N$ 即可。最优序列相互对应：
哈代常数由 $n^{-1/p}$ 逼近（第 21 题），
卡莱曼常数由 $n^{-1}$ 逼近 --- 两种情形中，调和型
数列恰好落在被平均的空间之外。

\textbf{25.} 对 $f = \mathbf 1_{\intcc01}$：$Hf(x) = 1$ 在
$\intoc01$ 上，$Hf(x) = \frac1x$ 对 $x > 1$，
$\norm{Hf}_2^2 = 1 + \int_1^\infty x^{-2}\dd x = 2$，于是
比值为 $\sqrt2 \approx 1.414$，约占尖锐界的 $71\%$。
对 $f_A$（$p = 2$）：$F(x) = 2(\sqrt x - 1)$ 在
$\intcc1A$ 上，故该范围内 $Hf_A = 2x^{-1/2} - 2x^{-1}$，
又对 $x > A$，$Hf_A(x) = 2(\sqrt A - 1)/x$。平方后
并积分，得
\begin{align*}
\int_1^A\bigl(2x^{-1/2} - 2x^{-1}\bigr)^2\dd x
&= 4\ln A - 12 + \frac{16}{\sqrt A} - \frac4A,\\
\int_A^{\infty}\frac{4(\sqrt A - 1)^2}{x^2}\,\dd x
&= 4 - \frac{8}{\sqrt A} + \frac4A,
\end{align*}
由此 $\norm{Hf_A}_2^2 = 4\ln A - 8 + 8A^{-1/2}$；两端
除以 $\norm{f_A}_2^2 = \ln A$，即得所述恒等式。取 $A
= \eu^{10}$ 时，$4 - \frac{8(1 - \eu^{-5})}{10} = 3.2054$，于是
比值为 $\sqrt{3.2054} \approx 1.790 < 2$。亏量
$4 - \norm{Hf_A}_2^2/\norm{f_A}_2^2 \sim 8/\ln A$ 只
以对数速率衰减：欲使比值达到 $1.99$，
需 $\ln A \approx 200$，即 $A \approx 10^{87}$。尖锐
常数是定理，而非实验。
\end{solution}
